\section{A proposed interface framework}
\label{sec:framework}

I use the five gaps to propose a common record for information passed between
components or robots.  The record states what is being sent, why the receiver
may use it, how to check it, and what to do if the check fails.  Formal claims
about refinement, realizability, controller synthesis, or the whole system
still need the appropriate behavioral models from \cref{sec:taxonomy}.

\begin{table}[htbp]
  \centering
  \caption{Interface-Contract Obligations Derived from the Five Gaps}
  \label{tab:obligations}
  \footnotesize
  \setlength{\tabcolsep}{3pt}
  \renewcommand{\arraystretch}{1.08}
  \begin{tabularx}{\textwidth}{@{}P{0.17\textwidth}P{0.33\textwidth}P{0.25\textwidth}Y@{}}
    \toprule
    Gap & Metadata and evidence carried by the interface & Required validators & Response on rejection \\
    \midrule
    G1: Semantic and physical validity
      & Canonical task and resource identifiers; hard constraints; frames and units; model, operating region, and margins
      & Registry resolution; task, feasibility, or certificate check
      & Clarify, reject, repair the task, or replan \\
    G2: Temporal validity
      & Event or forecast-target time; publication and availability times; validity interval; clock domain and error bound
      & Freshness, target-time, deadline, and clock-comparability predicates, e.g., \eqref{eq:temporal-admission}
      & Request fresh state, shorten the horizon, degrade, or stop \\
    G3: Uncertainty, provenance, and evidence
      & Uncertainty type, scope, correlations, calibration, and propagation rule; model and solver identity; lineage; evidence kind
      & Uncertainty-threshold, calibration, certificate, and provenance-policy checks
      & Re-estimate, re-solve, choose a robust alternative, or escalate \\
    G4: Implementation and lifecycle compatibility
      & Schema, protocol, implementation, dependency, and configuration identities; membership; resource ownership; team epoch
      & Supported-combination, epoch, ownership, and lifecycle-transition checks
      & Reconcile, downgrade, isolate, restart, or reallocate \\
    G5: Intervention, recovery, and attribution
      & Producer and consumer authority; safety overrides; acknowledgements; completion evidence; recovery owner
      & Authority, acknowledgement, intervention-report, and completion predicates
      & Transfer authority, coordinate replanning, stop, or request human action \\
    \bottomrule
  \end{tabularx}
\end{table}


\subsection{Interfaces as Producer--Consumer Edges}

I represent the deployed system as a directed graph whose vertices are
functions assigned to robots or services.  A \emph{vertical} edge
$e^{\mathrm v}=(r,k)\rightarrow(r,k+1)$ connects consecutive functions of robot
$r$, while a \emph{horizontal} edge $e^{\mathrm h}=(r,k)\rightarrow(s,l)$ with
$s\neq r$ connects robots or distributed services.  A horizontal edge is more
than another message path: its consumer may act on behalf of a different robot
and thereby change joint feasibility, resource ownership, or a peer's safety
margin.  Vertical edges usually share a clock, an implementation, and a
lifecycle; horizontal edges usually do not.

\subsection{Interface Contracts}

\begin{definition}[Interface contract]\label{def:interface-contract}
For an edge $e=(p,c)$ from producer $p$ to consumer $c$, an interface contract
is the tuple
\begin{equation}
  C_e=(P_e,M_e;\,A_e,G_e,V_e,E_e,R_e),
  \label{eq:contract}
\end{equation}
where $P_e$ is a typed payload; $M_e$ is typed metadata, such as identifiers,
frames and units, times, uncertainty semantics, implementation identities, team
epoch, and lineage; $A_e$ is the consumer's assumption; $G_e$ is the producer's
guarantee; $V_e$ is a set of validators or monitor references; $E_e$ is the
supporting evidence, such as a solver result, proof, calibration report, monitor
record, or test identifier; and $R_e$ is the response to rejection or unresolved
validity.
\end{definition}

A metadata field is contractual only when it appears in $A_e$, $G_e$, $V_e$, or
$R_e$.  \Cref{tab:obligations} maps each gap of \cref{sec:gaps} to the metadata,
validators, and recovery actions that an interface contract must specify; it is
the bridge from the review to the framework.  Provenance can reuse the
entity--activity--agent relations of W3C PROV \cite{w3c2013prov}.  The framework
keeps schema, protocol, and implementation identity separate because, as
\cref{sec:middleware-evidence} shows, syntactic compatibility does not establish
behavioral compatibility.

\subsection{Compatibility and Admission}
\label{sec:admission}

On a single edge, local compatibility requires the producer's guarantee to imply
the consumer's assumption:
\begin{equation}
  G_e\Longrightarrow A_e.
  \label{eq:compatibility}
\end{equation}
A consumer with several peer inputs must additionally check joint assumptions,
such as a common team epoch, consistent frames, unique resource ownership, and a
reconciliation rule after a partition; per-message checks cannot establish these
relations.  Condition \eqref{eq:compatibility} is therefore local.  Proving
that a composed team satisfies a global mission property requires defined
input--output ownership, feedback and temporal semantics, and a composition or
refinement calculus \cite{luckcuck2026compositional,mallozzi2023repair}.

Validators make the condition operational.  A syntactic validator checks
decoding, required fields, types, dimensions, units, frames, and schema
compatibility.  A semantic validator may call a solver, a runtime monitor, a
perception check, a calibrated statistical test, or a certificate verifier.
Each validator returns a value in $\{\true,\false,\unknown\}$.  Let
$\mathcal Q_e$ be the set of checks required on edge $e$, covering syntax,
time, uncertainty, compatibility, assumptions, guarantees, and provenance.
The edge admits the transfer at time $t$ if and only if every required check
returns \true:
\begin{equation}
  \operatorname{Admit}_e(t)\Leftrightarrow
  V_{e,q}(P_e,M_e,E_e,t)=\true\ \ \forall q\in\mathcal Q_e.
  \label{eq:admission}
\end{equation}
A required check that returns \unknown{} blocks admission, so missing metadata
or evidence leads to rejection rather than to silent acceptance; Kekatos
\emph{et al.}\ apply the same principle to missing monitoring evidence in robot
swarms \cite{kekatos2026swarmrta}.  Advisory checks outside $\mathcal Q_e$ may
return \unknown{} without blocking.  Condition \eqref{eq:admission} is what turns
descriptive fields into an enforced boundary: every required predicate takes
part in the execution decision.

\begin{table}[htbp]
  \centering
  \caption{Cross-Layer Failure-Attribution Taxonomy}
  \label{tab:attribution}
  \footnotesize
  \setlength{\tabcolsep}{3pt}
  \renewcommand{\arraystretch}{1.08}
  \begin{tabularx}{\textwidth}{@{}P{0.11\textwidth}P{0.15\textwidth}P{0.15\textwidth}YP{0.16\textwidth}P{0.15\textwidth}@{}}
    \toprule
    Origin & Local fault & Contract concern & First observation & Possible effect & Recovery authority \\
    \midrule
    Perception and estimation
      & Misassociation, bias, or clock error
      & Payload, frame, time, and uncertainty
      & Estimator innovation, semantic grounding, or controller residual
      & Wrong referent, localization drift, or unsafe state
      & Re-estimate, use a sensor fallback, or stop \\
    Semantic planning
      & Invented relation, capability, or priority
      & Intent, assumptions, and guarantees
      & Task-model validator or downstream feasibility check
      & Wrong or impossible assignment; repeated replanning
      & Clarify, repair, or replan \\
    Prediction
      & Biased or overconfident forecast
      & Horizon, uncertainty semantics, and validity evidence
      & Forecast residual or downstream feasibility check
      & Coordination failure or poor recovery
      & Re-estimate, shorten the horizon, or choose a robust plan \\
    Physical model
      & Missing constraint or inaccurate contact model
      & Feasibility claim and solver evidence
      & Optimizer status, controller residual, or contact event
      & Slip, saturation, collision, or damage
      & Re-optimize, switch to a fallback controller, or stop \\
    Runtime and network
      & Delay, loss, partition, version, or lifecycle fault
      & Age, clock, protocol, implementation identity, and assumptions
      & Consumer validator, missed update, exception, or teardown error
      & Inconsistent team state, abort, deadlock, or collision
      & Degrade, reconnect, reallocate, or restart \\
    Safety function
      & Local or unreported intervention
      & Authority, intervention evidence, and recovery
      & Safety monitor or upstream acknowledgement
      & Safe local motion but failed mission or new conflict
      & Coordinated replanning or human escalation \\
    Hardware
      & Sensor, actuator, energy, or compute fault
      & Capability and health guarantees
      & Health monitor, tracking residual, or peer timeout
      & Task abandonment or unsafe compensation by peers
      & Isolate the robot and reallocate its tasks \\
    \bottomrule
  \end{tabularx}
\end{table}


\subsection{Temporal Admission for Peer State}
\label{sec:temporal-admission}

The collision constraint of \eqref{eq:coupled-dynamics} shows which metadata a
horizontal interface must carry for a temporal validator to return \true.

\begin{proposition}[Temporal admission margin]
\label{prop:temporal-admission}
Let robots $r_i$ and $r_j$ have speeds bounded by $\bar v_i$ and $\bar v_j$, and
express all times in the consumer's clock.  Suppose the consumer holds position
estimates $\hat p_k$, $k\in\{i,j\}$, with stamped event times $\hat t_k$, such
that the true event time $t_k\le t$ satisfies $\lvert\hat t_k-t_k\rvert\le
\kappa_k$ and $\lVert\hat p_k-p_k(t_k)\rVert\le\varepsilon_k$.  For a decision
taken at time $t$ and held until $t+h$, let
$\bar\Delta_k=t+h-\hat t_k+\kappa_k$.  If
\begin{equation}
  \lVert\hat p_i-\hat p_j\rVert\ \ge\ d_{\min}
  +\sum_{k\in\{i,j\}}\bigl(\varepsilon_k+\bar v_k\bar\Delta_k\bigr),
  \label{eq:temporal-admission}
\end{equation}
then $\lVert p_i(s)-p_j(s)\rVert\ge d_{\min}$ for all $s\in[t,t+h]$.
\end{proposition}

\begin{proof}
Fix $s\in[t,t+h]$ and $k\in\{i,j\}$.  Because $t_k\le t\le s$ and
$t_k\ge\hat t_k-\kappa_k$, the true age satisfies
$0\le s-t_k\le\bar\Delta_k$.  The speed and error bounds give
$\lVert p_k(s)-\hat p_k\rVert\le\lVert p_k(s)-p_k(t_k)\rVert
+\lVert p_k(t_k)-\hat p_k\rVert\le\bar v_k\bar\Delta_k+\varepsilon_k$.  The
triangle inequality then yields $\lVert p_i(s)-p_j(s)\rVert\ge\lVert\hat
p_i-\hat p_j\rVert-\sum_k(\varepsilon_k+\bar v_k\bar\Delta_k)\ge d_{\min}$.
\end{proof}

Condition \eqref{eq:temporal-admission} is sufficient but not necessary, and it
is conservative because it ignores the direction of motion.  Its value here is
structural: every term on its right-hand side is a field that the interface must
carry or the consumer must already know.  The validator needs the event time
$\hat t_k$ rather than the reception time, the clock-error bound $\kappa_k$, an
error bound $\varepsilon_k$ with declared meaning, and the speed bound $\bar v_k$.
If any of these is missing, the validator returns \unknown{} and the decision is
not admitted.  Substituting the reception time for $\hat t_k$ understates the
age by the producer and network delay, which is exactly the failure described
in G2.

\begin{corollary}[Validity interval]\label{cor:validity-interval}
Under the assumptions of \cref{prop:temporal-admission}, let
$a_k=t-\hat t_k$ be the stamped age at decision time and
\[
  m(t)=\lVert\hat p_i-\hat p_j\rVert-d_{\min}
  -\sum_{k}\bigl(\varepsilon_k+\bar v_k(a_k+\kappa_k)\bigr)
\]
the separation margin that remains at that time.  Condition
\eqref{eq:temporal-admission} holds if and only if $h\le h^\star(t)$, where
\begin{equation}
  h^\star(t)=\frac{m(t)}{\bar v_i+\bar v_j}.
  \label{eq:validity-interval}
\end{equation}
\end{corollary}

\begin{proof}
Since $\bar\Delta_k=a_k+\kappa_k+h$, condition \eqref{eq:temporal-admission} is
equivalent to $m(t)\ge(\bar v_i+\bar v_j)\,h$, and $\bar v_i+\bar v_j>0$.
\end{proof}

If $h^\star(t)<0$, the pair cannot be admitted at time $t$ on these estimates.
Otherwise, the admitted state carries an explicit expiry time $t+h^\star(t)$,
after which the consumer must refresh it or invoke $R_e$.  The validity interval
that the gaps call for is thus computed by the consumer from fields carried on
the interface, not asserted by the producer.

\emph{Inference placement.}  \Cref{prop:temporal-admission} also prices the
choice between an on-board and a remote estimate of the same peer.  Suppose
robot $r_i$ can use either an on-board estimate of $p_j$, with error bound
$\varepsilon^{\mathrm L}$ and stamped age $a^{\mathrm L}=t-\hat t^{\mathrm L}$
on its own clock, or a remote estimate fused from more sensors, with error bound
$\varepsilon^{\mathrm R}$, stamped age $a^{\mathrm R}$, and clock-error bound
$\kappa^{\mathrm R}$.  The peer term in \eqref{eq:temporal-admission} is
$\varepsilon+\bar v_j(a+h+\kappa)$, so the on-board estimate consumes no more
of the separation margin than the remote one whenever
\begin{equation}
  \varepsilon^{\mathrm L}-\varepsilon^{\mathrm R}
  \ \le\ \bar v_j\bigl(a^{\mathrm R}+\kappa^{\mathrm R}-a^{\mathrm L}\bigr).
  \label{eq:placement}
\end{equation}
A remote estimate is worth its latency only if its accuracy advantage exceeds
the distance the peer can travel during the additional age and clock
uncertainty.  For example, with $\bar v_j=1.5$\,m/s, an additional
$120$\,ms of age, and a $10$\,ms clock-error bound, the remote estimate must be
at least $0.195$\,m more accurate.  The required advantage grows linearly with
peer speed, which is why on-board inference matters most for fast interactions,
whereas remote inference remains attractive for slow, wide-area decisions such
as task allocation \cite{kehoe2015cloud}.

\begin{table}[t]
  \centering
  \caption{Running command trace; all thresholds are scenario-specific contract parameters}
  \label{tab:command-trace}
  \scriptsize
  \begin{tabularx}{\textwidth}{@{}P{0.13\textwidth}P{0.27\textwidth}Y Y@{}}
    \toprule
    Function & Payload $P_e$ & Assumption, guarantee, and metadata & Validator and rejection $V_e,R_e$ \\
    \midrule
    Semantic task formation & Roles, sector identifier, inspection goal, and corridor-clear constraint & $A_e$: referents resolve and two capable robots exist; $G_e$: confirmed hard constraints remain explicit & Resolve against map and capability registries; clarify or reject an unresolved referent \\
    Predictive state & Occupancy and joint trajectories over horizon $H$ & $M_e$: event and target times, horizon, uncertainty semantics, and other-agent policy identity & Reject or shorten $H$ if age exceeds $\bar\Delta$ or uncertainty is not below its declared limit $\bar U$ \\
    Physical grounding & Reachable viewpoints and time-parameterized paths & $G_e$: separation is at least $d_{\min}$ and actuator limits hold; $E_e$: model, solver status, and minimum margins & Re-solve or return infeasible when solver status is unacceptable or a required margin is negative \\
    Distributed runtime & Assignments, paths, peer state, and acknowledgements & $M_e$: frame, units, schema/protocol/implementation identities, team epoch, ownership, and clock-error bound & Drop expired or wrong-epoch data; reconcile, degrade, or replan when age, clock, ownership, or compatibility checks fail \\
    Runtime safety & Admissible command or intervention record & $G_e$: conditional local-safety claim; $E_e$: active constraint, model/state/update assumptions, and intervention cause & Apply a feasible override and publish its cause; stop or escalate if the safety program is infeasible \\
    Physical robots & Execution state, health, resource release, and completion evidence & Capability and health guarantee; ownership and completion status & Isolate an unhealthy robot, reallocate its task, or request human action \\
    \bottomrule
  \end{tabularx}
\end{table}


\subsection{Recovery and Lineage}

A \false{} or a required \unknown{} invokes $R_e$: request clarification or
fresh data, re-estimate, replan, degrade, isolate, stop, or escalate.  The
rejection record names the failed predicate, the time, the evidence, and the
responsible interface.  A repaired result receives new lineage and does not
silently replace the rejected claim (\cref{fig:lifecycle}).

\begin{figure}[!t]
  \centering
  \begin{tikzpicture}[
    x=1cm,y=1cm,
    box/.style={draw,rounded corners=2pt,minimum height=7mm,align=center,
      font=\footnotesize,inner sep=2.5pt},
    arrow/.style={-{Latex[length=2mm]},line width=0.7pt},
    tag/.style={font=\scriptsize,align=center,inner sep=1pt}
  ]
    \definecolor{ContractBlue}{HTML}{DCEBFA}
    \definecolor{ValidatorGreen}{HTML}{DDEFD8}
    \definecolor{RecoveryOrange}{HTML}{F9E4C8}
    \definecolor{RejectRed}{HTML}{F7D8D5}

    \node[box,fill=ContractBlue,minimum width=36mm] (producer) at (0,0)
      {Producer output};
    \node[box,fill=ContractBlue,minimum width=52mm] (contract) at (0,-1.05)
      {Interface contract $C_e$: payload,\\metadata, assumptions, evidence};
    \node[box,fill=ValidatorGreen,minimum width=64mm] (validators) at (0,-2.35)
      {Required validators $V_{e,q}$, $q\in\mathcal Q_e$: syntax, time,\\
       uncertainty, compatibility, $A_e$, $G_e$, provenance};

    \coordinate (branch) at (0,-3.25);
    \node[box,fill=ValidatorGreen,text width=27mm] (admit) at (-2.05,-4.05)
      {Admit to consumer};
    \node[box,fill=RejectRed,text width=30mm] (reject) at (1.95,-4.05)
      {Reject; record failed predicate and evidence};
    \node[box,fill=RecoveryOrange,text width=30mm] (recover) at (1.95,-5.25)
      {Invoke $R_e$: refresh, re-estimate, replan, degrade, or stop};
    \node[box,fill=ContractBlue,text width=30mm] (repair) at (1.95,-6.40)
      {New contract instance with new lineage};

    \draw[arrow] (producer) -- (contract);
    \draw[arrow] (contract) -- (validators);
    \draw[line width=0.7pt] (validators.south) -- (branch);
    \draw[arrow] (branch) -| (admit.north);
    \draw[arrow] (branch) -| (reject.north);
    \node[tag,anchor=south] at (-2.05,-3.22) {all required $=\true$};
    \node[tag,anchor=south] at (1.95,-3.22) {\false{} or required \unknown};
    \draw[arrow] (reject) -- (recover);
    \draw[arrow] (recover) -- (repair);
    \draw[arrow] (repair.east) -- ++(0.35,0) |- (contract.east);
  \end{tikzpicture}
  \caption{Admission and recovery on one interface.  A required check that
  returns \unknown{} blocks admission like \false.  A successful repair creates a
  new contract instance and provenance record while preserving the rejected
  claim.}
  \label{fig:lifecycle}
\end{figure}


\subsection{Failure Attribution}
\label{sec:attribution}

The same physical outcome can have several origins.  A collision may begin with
an incorrect semantic relation, a biased prediction, a missing physical
constraint, stale peer state, or a hardware fault.  Conversely, downstream
checks can turn the same originating fault into a rejection, a degraded mission,
or a safe stop.  The taxonomy in \cref{tab:attribution} therefore records the
origin, the contract concern, the first observation, the effect, and the
recovery authority as separate fields; the first observation is the earliest
interface at which evidence of the fault can become available.

\subsection{Worked Command Trace}
\label{sec:trace}

Consider the instruction \emph{``Use two robots to inspect sector~B while
keeping the corridor clear.''}  \Cref{tab:command-trace} follows the command
from resolved roles and constraints to trajectories, runtime state, safety
intervention, and completion evidence.  Each row lists the information that the
next validator requires in addition to the nominal output.  The trace is a
reference representation rather than a universal ROS message definition;
thresholds and recovery policies remain scenario-specific.
